\documentclass[11pt,twocolumn]{article}

\usepackage[a4paper,margin=0.75in,columnsep=0.3in]{geometry}
\usepackage[T1]{fontenc}
\usepackage{amsmath,amssymb,amsthm,mathtools}
\usepackage{bm}
\usepackage{graphicx}
\usepackage{booktabs}
\usepackage{array}
\usepackage{enumitem}
\usepackage{float}
\usepackage[numbers,sort&compress]{natbib}
\usepackage{hyperref}
\hypersetup{colorlinks=true,linkcolor=blue,citecolor=blue,urlcolor=blue}

\newtheorem{theorem}{Theorem}
\newtheorem{definition}{Definition}
\newtheorem{postulate}{Postulate}
\newtheorem{corollary}{Corollary}
\newtheorem{proposition}{Proposition}

\title{Unified Dynamics:\\Emergent Spacetime, Gauge Structure, Matter, Cosmology,\\and Dimensional Ontology from a Single Tetrahedral Primitive}
\author{%
Benjamin Walker Mayes\thanks{*Contact author: ben.w.mayes@gmail.com}\\[4pt]
{\small Unaffiliated, Cedar Park, Texas 78613, USA}\\[8pt]
Arthur Eugene Mayes\\
{\small Loyola University School of Law, USA}\\[4pt]
Dustin Huff,\quad Brayden Ross Sanders,\quad Jeffrey A.\ Dewey\\
{\small Unaffiliated, USA}\\[4pt]
Brandon Bevin\\
{\small Georgia Southern University, USA}\\[4pt]
A.\ Christopher Beckingham\\
{\small Queen's University, Canada}\\[4pt]
Jonathan Engle\\
{\small Unaffiliated, USA}%
}
\date{}

\begin{document}
\maketitle

\begin{abstract}
We present Unified Dynamics (UD), a framework in which spacetime geometry, gauge interactions, matter, quantum measurement, dark matter, baryogenesis, the cosmological constant, black hole information, and the dimensionality of space all emerge from admissible transport on a single combinatorial primitive: the tetrahedron~$K_4$.
An iterated function system (IFS) with contraction ratio $r=4^{-1/3}$ generates a fractal lattice whose combinatorial constants---overlap deficit $\Delta=7$, state-space dimension $N=192$, master scalar $x=7/192$---determine all structural parameters without fitting.

The gauge group $\mathrm{SU}(3)\times\mathrm{SU}(2)\times\mathrm{U}(1)$ emerges from $S_4$ representation theory with $\sin^2\theta_W=3/11$.
A Yang--Mills mass gap $\Delta_{\mathrm{mass}}\geq 3/\pi$ is proved, implying confinement with string tension $\sigma=9/(4\pi^2 a_0)$.
Electroweak symmetry breaking is driven by the $\mathbb{Z}_2$ ferromagnetic vacuum.
A nonlinear fermion equation with derived coupling $\lambda_{\mathrm{NL}}=441/(2048\pi^2)$ produces three generations; the lepton mass tower matches experiment to $0.2\%$.
The Born rule emerges from carrier-viability selection; entanglement from $\mathbb{Z}_2$ carrier conservation.
The continuum Einstein--Hilbert action is the unique Banach fixed point of the IFS coarse-graining.

A two-stage electroweak transition produces dark-matter relics (gauge-singlet composites, $\Omega_{\mathrm{DM}}\approx 43/228\approx 19\%$), baryogenesis from CP-violating framing dynamics ($\eta_B\sim 10^{-5}$ pre-washout), and a structurally screened cosmological constant ($\rho_\Lambda\propto x^{p}$).
Black hole information is preserved by carrier conservation.

The spectral dimension of the Sierpi\'nski tetrahedron is $d_s=4/3$.
We prove that $d=3$ is the unique stable closure dimension: for $d\geq 4$ the binding coupling falls below threshold and stable structures require composites of $d=3$ closures, with a coherence ceiling at $d_{\max}\approx 7$.
Neutrino masses near $35$~meV and PMNS mixing angles are derived from generation-spacing mismatch.

The theory has one input scale ($\nu_0\approx 83$~keV) and zero free parameters.
\end{abstract}

%═══════════════════════════════════════════════════════════
% PART I — FOUNDATIONS
%═══════════════════════════════════════════════════════════

\section{Introduction}
\label{sec:intro}

A unified description of nature should derive geometry, gauge interactions, matter, quantum mechanics, cosmology, and dimensionality from a single starting point.
General relativity~\cite{einstein1915} treats geometry as dynamical; quantum field theory~\cite{yangmills1954} assumes a fixed arena.
Unified Dynamics (UD) starts earlier: physical reality is a discrete packet-routing substrate whose admissible transport, closure, and selective retention generate all known physics as derived consequences.

The framework rests on one primitive: the tetrahedron $K_4$, the unique self-dual minimal 3D closure.
An iterated function system~\cite{hutchinson1981,falconer2003} generates a fractal lattice whose combinatorial constants determine every structural parameter.
The present paper consolidates the full derivation chain---from $K_4$ through gauge structure, symmetry breaking, matter, quantum measurement, cosmology, dark matter, baryogenesis, black hole information, the dimensionality of space, and the mass spectrum---with explicit equations, zero free parameters, and falsifiable predictions.

\section{Foundational Axioms}
\label{sec:axioms}

\begin{postulate}[Packet substrate]
Each site $i$ on the tetrahedral IFS lattice carries state variables
\begin{equation}
X_i = (\rho_i,\,\theta_i,\,\chi_i,\,\sigma_i,\,b_i),
\label{eq:packet}
\end{equation}
where $\rho_i$ is burden density, $\theta_i\in[0,2\pi)$ is transport phase, $\chi_i\in\{0,1\}$ is closure state, $\sigma_i\in\{+1,-1\}$ is framing, and $b_i\in[0,1]$ is boundary integrity.
\end{postulate}

\begin{postulate}[Transport]
Neighboring packets are connected by transport operators $U_{ij}$; composition around a closed loop $C$ defines holonomy $U_C=\prod_{(ij)\in C}U_{ij}$~\cite{wilson1974}.
\end{postulate}

\begin{postulate}[Admissibility]
The physical state space is $\mathcal{T}_{\mathrm{phys}}=\{\Psi\mid\mathcal{A}_{\mathrm{closure}}(\Psi)=0\}$.
\end{postulate}

Six axioms constrain the grammar:
\textbf{(A1)~Seed}---persistent families begin from minimal closure;
\textbf{(A2)~Threshold}---new regimes cost at least~$S_0$;
\textbf{(A3)~Lift}---higher structures from iterated transport;
\textbf{(A4)~Orientation}---$\mathbb{Z}_2$ framing determines statistics;
\textbf{(A5)~Exchange}---transport modes define interactions;
\textbf{(A6)~Dressing}---observables are spines plus corrections.

\section{Tetrahedral Combinatorics and IFS}
\label{sec:K4}

$K_4$ has $V{=}4$, $E{=}6$, $F{=}4$, $\chi{=}2$, $|S_4|{=}24$.
Four contraction maps $\{f_k\}$ with ratio $r=4^{-1/3}$ generate the Sierpi\'nski tetrahedron~\cite{hutchinson1981}.
\begin{align}
N &= |S_4|\times 2^3 = 192,\label{eq:N}\\
\Delta &= \tbinom{4}{2}+1 = 7,\quad x = 7/192,\label{eq:x}\\
\eta_R &= 185/192.\label{eq:eta}
\end{align}

\section{Packet Action and Transport}
\label{sec:action}

The packet action is
\begin{multline}
S_{\mathrm{pkt}} = \sum_C \kappa_C\,\mathrm{Tr}(I{-}U_C) + \sum_i \lambda_B(\rho_i{-}\rho_0)^2\\
+ \sum_{\langle ij\rangle}\eta_O(\theta_i{-}\theta_j)^2,
\label{eq:Spkt}
\end{multline}
with $\kappa_C=3/(4\pi)$, $\lambda_B=\kappa_C S_0$, $\eta_O=1/2$.
The transport seed factorizes as $T_0=(1/2)(7/16)(1/6)=7/192$, and the path action selects the least-strain coherent path.

%═══════════════════════════════════════════════════════════
% PART II — GAUGE STRUCTURE AND GEOMETRY
%═══════════════════════════════════════════════════════════

\section{Structural Constants}
\label{sec:constants}

\begin{theorem}[Thresholds and coupling]
\label{thm:constants}
$S_0=(E\Delta{+}1)/\Delta=43/7$; $H=2\cdot 2^3{+}1{-}1/6=101/6$; $\alpha_0=\Delta/(5N)=7/960$; $\kappa_C=3/(4\pi)$ from discrete Gauss--Bonnet~\cite{hirani2003}.
\end{theorem}

\section{Gauge Group}
\label{sec:gauge}

\begin{theorem}
$\mathrm{SU}(3){\times}\mathrm{SU}(2){\times}\mathrm{U}(1)$ from $S_4$ on $K_4$: $g_3^2{=}3/\pi$ (faces), $g_2^2{=}3/(2\pi)$ (edge pairs), $g_1^2{=}9/(16\pi)$ (vertices).
Ratios $48{:}24{:}9$ give $\sin^2\theta_W=3/11$.
\end{theorem}

\section{Electroweak Symmetry Breaking}
\label{sec:ewsb}

\begin{theorem}
The $\mathbb{Z}_2$ ferromagnetic vacuum $\sigma_i{=}{+}1$ breaks $\mathrm{SU}(2)\to\mathrm{U}(1)_{\tau_3}$.
The Higgs is the framing-alignment amplitude mode.
\end{theorem}

\section{Mass Gap and Confinement}
\label{sec:massgap}

\begin{theorem}[Mass gap]
$\Delta_{\mathrm{mass}}\geq\kappa_C\Delta_{\mathrm{curv}}=3/\pi>0$.
The gap survives the continuum: IFS growth $\eta_R/r\approx 1.53>1$.
Discrete curvature follows Regge~\cite{regge1961}.
\end{theorem}

\begin{theorem}[Confinement]
The mass gap implies an area-law Wilson loop~\cite{wilson1974} with string tension $\sigma=\kappa_C\Delta_{\mathrm{mass}}/a_0=9/(4\pi^2 a_0)$.
\end{theorem}

\section{Emergent Geometry}
\label{sec:geometry}

\begin{theorem}[Continuum limit]
$S_{\mathrm{EH}}=(16\pi G)^{-1}\!\int d^4x\sqrt{g}\,R$ is the unique Banach fixed point of IFS coarse-graining, with $G=\kappa_C a_0^2$.
The vierbein from DEC face normals~\cite{hirani2003} yields Einstein equations~\cite{einstein1915}; the accommodation architecture organizes departures~\cite{mayes2026grav}.
\end{theorem}

%═══════════════════════════════════════════════════════════
% PART III — MATTER AND QUANTUM
%═══════════════════════════════════════════════════════════

\section{Fermion Equation and Cross-Coupling}
\label{sec:fermion}

\begin{theorem}[Cross-coupling]
$g_{\chi\sigma}=\Delta(\Delta{-}1)/(2^{E+1}\pi)=21/(64\pi)$, where $21=\binom{7}{2}$ counts pairwise deficit modes, $64=2^E$, and $\pi$ is the consumed-loop phase.
\end{theorem}

\begin{theorem}[Nonlinear fermion equation]
$(i\gamma^\mu D_\mu-m)\psi=\lambda_{\mathrm{NL}}|\bar\psi\psi|\psi$, with $\lambda_{\mathrm{NL}}=g_{\chi\sigma}^2/J_\sigma=441/(2048\pi^2)$.
The $\gamma$-matrices emerge from $\deg(v)=3$.
\end{theorem}

\section{Mass Spectrum}
\label{sec:mass}

The mass hierarchy: $m_\tau/m_\mu=H=101/6\approx 16.83$ (measured $16.82$, $0.09\%$).
Three generations from $b$-channel budget $<H^2$.

The muon mass with sub-leading IFS corrections:
\begin{equation}
\frac{m_\mu}{m_e} = 4^4\!\left(\frac{185}{192}\right)^{\!6}\!\left(1+\frac{\Delta{+}C^*}{NE}+\frac{F}{NEd}\right)\approx 206.3,
\label{eq:muon}
\end{equation}
matching the measured $206.8$ to $0.2\%$.
Nuclear binding follows Weizs\"acker~\cite{weizsacker1935} with UD coefficients.

Neutrino masses from second-order $x$ coupling with consumed-loop suppression:
\begin{equation}
m_{\nu_1} = \nu_0 x^2/\pi \approx 35\text{ meV},
\label{eq:neutrino}
\end{equation}
consistent with cosmological bounds~\cite{planck2018}.
PMNS angles from generation-spacing mismatch ($12{:}6$ vs $1{:}1$):
$\theta_{12}=\arctan(1/2)\approx 27^\circ$, $\theta_{23}=45^\circ$, $\theta_{13}\approx 6.6^\circ$.

\section{Quantum Measurement and Entanglement}
\label{sec:quantum}

The carrier $\mathcal{C}_i=(C_\rho,C_\theta,C_\chi,C_\sigma,C_b)$ has viability $\mathcal{V}\geq 1$.
Five measurement laws without a measurement axiom:
branch generation ($\eta_B=7/16$),
Born rule from Gleason~\cite{gleason1957},
decoherence ($\eta_O N_{\mathrm{env}}$),
context stability ($e^{-\nu_0 N}$),
frequency bridge.

\begin{theorem}[Entanglement]
$\mathbb{Z}_2$ carrier conservation at splitting vertices ($\sigma_A\sigma_B=\sigma_{\mathrm{parent}}$) produces entanglement.
Bell violation $S=2\sqrt{2}$ from Born-rule weighting.
Cost: $\nu_0/\pi$ per pair.
The Madelung transform~\cite{madelung1927} confirms the Schr\"odinger limit.
\end{theorem}

%═══════════════════════════════════════════════════════════
% PART IV — COSMOLOGY
%═══════════════════════════════════════════════════════════

\section{Framing Phase Transition}
\label{sec:transition}

The framing sector is a ferromagnetic Ising model with $J_\sigma=1/2$, $z=6$, critical temperature $T_c=3\nu_0$.
The electroweak process has two stages:
\textbf{Stage~1} ($T\sim v_{\mathrm{EW}}$): $\mathrm{SU}(2)$ breaking, baryogenesis (sphalerons active).
\textbf{Stage~2} ($T\sim 3\nu_0$): framing Ising ordering, dark-matter freeze-out, trapped defects.

\section{Dark Matter}
\label{sec:dm}

\begin{theorem}
Stable neutral composites---4-vertex $K_4$ closures with $\sum_i\sigma_i=0$---are gauge singlets, closure-protected, gravitationally active.
Mass $m_{\mathrm{DM}}=4\nu_0 S_0\approx 2$~MeV.
\end{theorem}

\begin{proposition}[Relic fraction]
$\Omega_{\mathrm{DM}}/\Omega_{\mathrm{total}}=xS_0/(xS_0+1{-}x)=43/228\approx 0.189$.
Measured: $0.27$ ($30\%$ discrepancy, may close with dressing).
\end{proposition}

\section{Baryogenesis}
\label{sec:baryo}

\begin{proposition}
CP-violating framing cross-coupling contributes to baryogenesis during Stage~1~\cite{sakharov1967,kuzmin1985}.
Pre-washout $\eta_B\sim 10^{-5}$; achieving $6\times 10^{-10}$ requires washout $\sim 10^5$.
\end{proposition}

\section{Cosmological Constant}
\label{sec:cc}

Trapped framing defects with energy $\Delta E=6\nu_0$ and IFS screening give $\rho_\Lambda=A\nu_0^4 x^{2+n_{\mathrm{screen}}}$.
The screening exponent $n_{\mathrm{screen}}\approx 19$ (cosmological IFS sweep) yields $\rho_\Lambda\sim 10^{-123}$ in Planck units, within $1.3$ orders of the observed value.

\section{Black Hole Information}
\label{sec:bh}

\begin{theorem}
Carrier conservation through all local transport implies no information loss.
The horizon is effective, not fundamental.
Entangled carrier pairs leak correlations through the horizon~\cite{hawking1975}.
Bekenstein--Hawking entropy~\cite{bekenstein1973}: $S_{\mathrm{BH}}=A/(4\kappa_C a_0^2)$.
Page time~\cite{page1993}: information released as $S_{\mathrm{BH}}$ drops to half.
\end{theorem}

%═══════════════════════════════════════════════════════════
% PART V — DIMENSIONAL ONTOLOGY
%═══════════════════════════════════════════════════════════

\section{Why Three Dimensions}
\label{sec:d3}

The $d$-simplex $K_{d+1}$ has $N(d)=(d{+}1)!2^d$, $\Delta(d)=d(d{+}1)/2{+}1$, $\lambda_{\mathrm{crit}}(d)=d/N(d)$.

\begin{theorem}[$D=3$ uniqueness]
$d=3$ is the unique dimension satisfying: (i)~$F_2=V$, (ii)~$x\in(0.01,0.1)$, (iii)~$\lambda_{\mathrm{NL}}>d/N$, (iv)~$\eta_R\in(0.9,0.99)$.
\end{theorem}

\section{Higher-Dimensional Growth}
\label{sec:growth}

For $d\geq 4$, $\lambda_{\mathrm{NL}}<d/N$: lone primitives cannot bind.
Composites of $n_c$ tetrahedra have $\lambda^{\mathrm{comp}}=n_c^2\lambda_{\mathrm{NL}}(d)$.
Three classes: seed ($n_c{=}4$), stable composite ($n_c{=}8$), shell molecule ($n_c{\geq}16$).
The coherence ceiling $n_c^{\max}=\lfloor 185/3\rfloor=61$ caps the ladder at $d_{\max}\approx 7$.
Rupture produces $d{=}3$ remnants.

%═══════════════════════════════════════════════════════════
% PART VI — SPECTRAL GEOMETRY AND RG
%═══════════════════════════════════════════════════════════

\section{Spectral Dimension}
\label{sec:spectral}

\begin{theorem}
$d_s=2\ln 4/\ln 8=4/3$ from decimation on the Sierpi\'nski tetrahedron~\cite{rammal1984}.
The RG correction factor is $F_{\mathrm{IFS}}=d_s/4=1/3$.
\end{theorem}

\section{IFS Level Hierarchy}
\label{sec:hierarchy}

The lattice spacing $a_0=\sqrt{\kappa_C}/M_{\mathrm{Pl}}$ places the IFS at the Planck scale.
Level~0 (Planck) $\to$ ${\sim}85$ (electroweak) $\to$ ${\sim}102$ ($\nu_0$) $\to$ ${\sim}130$ (QCD) $\to$ ${\sim}190$ (Hubble).
IFS beta functions use $b_i^{\mathrm{IFS}}=b_i^{\mathrm{SM}}\cdot F_{\mathrm{IFS}}\cdot(\text{thresholds})$; the remaining $5$--$8\times$ $\alpha_s$ discrepancy awaits the carrier excitation spectrum.

%═══════════════════════════════════════════════════════════
% PART VII — PREDICTIONS AND CLOSURE
%═══════════════════════════════════════════════════════════

\section{Unified Predictions}
\label{sec:predictions}

Table~\ref{tab:unified} collects all zero-parameter predictions across gauge structure, matter, cosmology, and dimensionality.
Every entry is derived from $K_4$ combinatorics plus the single input scale $\nu_0\approx 83$~keV.

\begin{table}[h!]
\centering
\caption{Zero-parameter predictions from $K_4$.}
\label{tab:unified}
\small
\begin{tabular}{@{}l@{\hskip 4pt}l@{\hskip 4pt}l@{\hskip 4pt}l@{}}
\toprule
Quantity & Prediction & Meas.\ & Acc.\ \\
\midrule
\multicolumn{4}{@{}l}{\textit{Gauge \& geometry}}\\
$\alpha_0^{-1}$ & $960/7$ & $137.04$ & 0.08\% \\
$\sin^2\theta_W$ & $3/11$ & $0.231$ & unif.\ \\
$g_3^2{:}g_2^2{:}g_1^2$ & $48{:}24{:}9$ & --- & lattice \\
$\Delta_{\mathrm{mass}}$ & ${\geq}3/\pi$ & ${>}0$ & struct.\ \\
$\sigma$ & $9/(4\pi^2 a_0)$ & --- & struct.\ \\
$\kappa_C$ & $3/(4\pi)$ & $0.24$ & 0.5\% \\
$G$ & $\kappa_C a_0^2$ & --- & lattice \\
\midrule
\multicolumn{4}{@{}l}{\textit{Matter}}\\
$m_\mu/m_e$ & $206.3$ & $206.8$ & 0.2\% \\
$m_\tau/m_\mu$ & $101/6$ & $16.82$ & 0.09\% \\
$N_{\mathrm{gen}}$ & $3$ & $3$ & exact \\
$g_{\chi\sigma}$ & $21/(64\pi)$ & --- & struct.\ \\
$m_{\nu_1}$ & $35$ meV & $<120$ & consist.\ \\
$\theta_{23}$ & $45^\circ$ & $42$--$50^\circ$ & consist.\ \\
$\theta_{13}$ & $6.6^\circ$ & $8.6^\circ$ & 23\% \\
\midrule
\multicolumn{4}{@{}l}{\textit{Cosmology}}\\
$n_s$ & $185/192$ & $0.965$ & Planck \\
$r_T$ & $7/192$ & ${<}0.036$ & consist.\ \\
$\Omega_{\mathrm{DM}}$ & $43/228$ & $0.27$ & 30\% \\
No info loss & carrier & --- & struct.\ \\
\midrule
\multicolumn{4}{@{}l}{\textit{Dimensionality}}\\
$d_{\mathrm{stable}}$ & $3$ & $3$ & exact \\
$d_s$ & $4/3$ & --- & struct.\ \\
$d_{\max}$ & $\approx 7$ & --- & predict.\ \\
\bottomrule
\end{tabular}
\end{table}

\section{Computational Model}
\label{sec:computation}

The simulation kernel $H_{\min}=H_\theta+H_\chi+H_\sigma+H_{\mathrm{int}}$ has all couplings derived; proxy runs at IFS levels $n=5$--$7$ yield $d_f\approx 1.53$.

\section{Falsifiers}
\label{sec:falsifiers}

(F1)~large non-GR weak-field deviations;
(F2)~no controlled GR limit;
(F3)~$n_s\notin[0.960,0.968]$;
(F4)~monopoles discovered;
(F5)~$\alpha_0^{-1}\neq 960/7$;
(F6)~no stable packets in simulation;
(F7)~$r_T>0.048$;
(F8)~DM with gauge interactions;
(F9)~baryogenesis incompatible with $\mathbb{Z}_2$ framing;
(F10)~information loss in BH evaporation;
(F11)~$\Lambda$ exactly zero;
(F12)~fourth fermion generation;
(F13)~stable lone $d{\geq}4$ bound states;
(F14)~$d_s\neq 4/3$ on tetrahedral fractal.

\section{Discussion}
\label{sec:discussion}

UD shares goals with lattice gauge theory~\cite{wilson1974}, loop quantum gravity~\cite{rovelli2004}, causal dynamical triangulations, and string theory~\cite{polyakov1987}, but differs in deriving ALL constants from a single combinatorial primitive.
The answer to ``why three dimensions?''\ does not invoke anthropic reasoning~\cite{weinberg1987} or compactification; it follows from the binding criterion $\lambda_{\mathrm{NL}}>d/N$.
The entropy bound~\cite{bekenstein1973} is automatically respected.

Open computations: full RG flow with IFS beta functions, the carrier excitation spectrum from $H_{\min}$ simulation, precise relic abundance, baryogenesis washout, and the Jaffe--Witten spectral verification.
These are computations within the specified framework, not structural gaps.

\section{Conclusion}
\label{sec:conclusion}

A single tetrahedral primitive generates: the Standard Model gauge group, electroweak symmetry breaking, a mass gap with confinement, the Einstein--Hilbert action, three generations of fermions with the observed mass hierarchy, quantum measurement and entanglement, dark matter, baryogenesis, a structurally screened cosmological constant, information-preserving black holes, the unique stability of three spatial dimensions, neutrino masses, and the spectral dimension of the substrate.
The theory has one input scale and zero free parameters.
Every structural constant traces to $K_4$ combinatorics.

The strongest falsifier is F6: if no stable closure-capable packets emerge in simulation, the substrate ontology fails.
A positive signal would provide the first evidence that observable physics is generated by a transport--closure architecture.

\section*{Acknowledgments}
The authors thank iterative computational dialog and theory-development workflows for enabling the progressive consolidation of the Unified Dynamics framework.

\section*{Author contributions}
B.W.M.\ conceived the framework, derived all results, and wrote the manuscript.
A.E.M., D.H., B.R.S., J.A.D., B.B., A.C.B., and J.E.\ contributed through discussion, review, and validation.

\smallskip
\noindent\textbf{Data availability statement.}
This paper has no associated data.
This is a theoretical study; no experimental data were generated or analyzed.

\begin{thebibliography}{99}
\bibitem{einstein1915} A.~Einstein, ``Die Feldgleichungen der Gravitation,'' Sitzungsber.\ Preuss.\ Akad.\ Wiss., 844--847 (1915).
\bibitem{yangmills1954} C.~N.~Yang and R.~L.~Mills, ``Conservation of isotopic spin and isotopic gauge invariance,'' Phys.\ Rev.\ \textbf{96}, 191--195 (1954).
\bibitem{wilson1974} K.~G.~Wilson, ``Confinement of quarks,'' Phys.\ Rev.\ D \textbf{10}, 2445--2459 (1974).
\bibitem{rovelli2004} C.~Rovelli, \textit{Quantum Gravity} (Cambridge University Press, Cambridge, 2004).
\bibitem{hutchinson1981} J.~E.~Hutchinson, ``Fractals and self-similarity,'' Indiana Univ.\ Math.\ J.\ \textbf{30}, 713--747 (1981).
\bibitem{falconer2003} K.~Falconer, \textit{Fractal Geometry}, 2nd ed.\ (Wiley, Chichester, 2003).
\bibitem{hirani2003} A.~N.~Hirani, ``Discrete exterior calculus,'' Ph.D.\ thesis, California Institute of Technology (2003).
\bibitem{gleason1957} A.~M.~Gleason, ``Measures on the closed subspaces of a Hilbert space,'' J.\ Math.\ Mech.\ \textbf{6}, 885--893 (1957).
\bibitem{madelung1927} E.~Madelung, ``Quantentheorie in hydrodynamischer Form,'' Z.\ Phys.\ \textbf{40}, 322--326 (1927).
\bibitem{planck2018} Planck Collaboration, ``Planck 2018 results.\ VI.\ Cosmological parameters,'' Astron.\ Astrophys.\ \textbf{641}, A6 (2020).
\bibitem{ligo2016} LIGO Scientific and Virgo Collaborations, ``Observation of gravitational waves from a binary black hole merger,'' Phys.\ Rev.\ Lett.\ \textbf{116}, 061102 (2016).
\bibitem{weizsacker1935} C.~F.~von Weizs\"acker, ``Zur Theorie der Kernmassen,'' Z.\ Phys.\ \textbf{96}, 431--458 (1935).
\bibitem{polyakov1987} A.~M.~Polyakov, \textit{Gauge Fields and Strings} (Harwood Academic, Chur, 1987).
\bibitem{regge1961} T.~Regge, ``General relativity without coordinates,'' Nuovo Cimento \textbf{19}, 558--571 (1961).
\bibitem{rammal1984} R.~Rammal and G.~Toulouse, ``Random walks on fractal structures and percolation clusters,'' J.\ Physique Lett.\ \textbf{44}, L13--L22 (1983).
\bibitem{sakharov1967} A.~D.~Sakharov, ``Violation of CP invariance, C asymmetry, and baryon asymmetry of the universe,'' JETP Lett.\ \textbf{5}, 24--27 (1967).
\bibitem{kuzmin1985} V.~A.~Kuzmin, V.~A.~Rubakov, and M.~E.~Shaposhnikov, ``On anomalous electroweak baryon-number non-conservation in the early universe,'' Phys.\ Lett.\ B \textbf{155}, 36--42 (1985).
\bibitem{hawking1975} S.~W.~Hawking, ``Particle creation by black holes,'' Commun.\ Math.\ Phys.\ \textbf{43}, 199--220 (1975).
\bibitem{bekenstein1973} J.~D.~Bekenstein, ``Black holes and entropy,'' Phys.\ Rev.\ D \textbf{7}, 2333--2346 (1973).
\bibitem{page1993} D.~N.~Page, ``Information in black hole radiation,'' Phys.\ Rev.\ Lett.\ \textbf{71}, 3743--3746 (1993).
\bibitem{weinberg1987} S.~Weinberg, ``Anthropic bound on the cosmological constant,'' Phys.\ Rev.\ Lett.\ \textbf{59}, 2607--2610 (1987).
\bibitem{mayes2026grav} B.~W.~Mayes, ``Unified Dynamics as an accommodation-channel extension of general relativity,'' Zenodo (2026), \href{https://doi.org/10.5281/zenodo.19582493}{doi:10.5281/zenodo.19582493}.
\end{thebibliography}

\end{document}
